\section{Die Diagnose}

\textbf{[K]} Aus A085 (die Naturkonstanten $c, \hbar, k_B$) und A130 ($\alpha$)
lässt sich ein einziger, scharfer Satz ziehen, der bisher nur verstreut
anklang:
\begin{quote}
\emph{Jede Konstante, die sich durch eine widerspruchsfreie Einheiten- oder
Normierungswahl auf Eins setzen lässt, ist damit als \emph{nicht grundlegend}
erwiesen. Sie ist ein Umrechnungsfaktor, kein Freiheitsgrad.}
\end{quote}

Die Setzbarkeit auf Eins ist also nicht bloß eine bequeme Schreibweise, sondern
ein \emph{Test}. Was auf Eins gebracht werden kann, ohne dass eine physikalische
Aussage verloren geht, war nie ein Inhalt der Natur, sondern eine Eigenschaft
des benutzten Maßsystems. Der berühmte Zahlenwert -- $c = 299\,792\,458$,
$1/\alpha = 137{,}036$ -- misst dann die \emph{historische Einheit}, nicht die
Natur.

\section{Zwei Wege auf Eins, ein Ergebnis}

\textbf{[K]} Es gibt zwei verschiedene Gründe, warum eine Konstante auf Eins
gebracht werden kann, und es ist wichtig, sie zu unterscheiden -- weil das
Ergebnis in beiden Fällen dasselbe ist, der Grund aber nicht.

\textbf{[B]} \emph{Erstens über eine physikalische Identität.} $c$, $\hbar$,
$k_B$ werden Eins, weil die Zeit-Masse-Dualität ihre Dimensionen
\emph{identifiziert} (A085): Masse, Energie, Temperatur und inverse Zeit sind
eine Größe in vier Kleidungen. Die $=1$ ist hier das Ablesen einer Identität.
Der Grund ist inhaltlich.

\textbf{[K]} \emph{Zweitens über eine Normierung.} $\alpha$ wird Eins, weil man
die Ladung so normiert, dass $e^2 = 4\pi$ (A130). Der Grund ist eine
Einheitenwahl, kein physikalischer Satz.

\textbf{[K]} Der springende Punkt: \emph{für die Diagnose ist der Unterschied
gleichgültig}. Ob eine Konstante über eine Identität oder über eine Normierung
auf Eins fällt -- in beiden Fällen zeigt die Setzbarkeit, dass sie kein eigener
Freiheitsgrad ist. Der Grund entscheidet, \emph{wie tief} die Buchhaltung
verankert ist (bei der Dualität in einer physikalischen Identität, bei
$\alpha$ in einer Konvention), nicht \emph{ob} die Konstante grundlegend ist.
Grundlegend ist keine von beiden.

\section{Was sich \emph{nicht} auf Eins setzen lässt}

\textbf{[K]} Der Satz ist nur scharf, wenn seine Umkehrung stimmt: es muss
Größen geben, die sich \emph{durch keine} Einheitenwahl auf Eins bringen lassen
-- und genau die sind der Inhalt.

\textbf{[B]} Zwei Arten überleben den Test:
\begin{itemize}
  \item \textbf{Der geometrische Parameter $\xi$.} Er ist ein reines
    Zahlenverhältnis der Packungsgeometrie (A020). Kein Wechsel der Meter, der
    Sekunde oder der Ladungseinheit ändert ihn; er ist $4/30000$ in jedem
    Maßsystem. $\xi$ lässt sich nicht auf Eins normieren, ohne die Geometrie
    selbst zu ändern.
  \item \textbf{Dimensionslose Verhältnisse.} $m_\mu/m_e = 206{,}77$ ist der
    Quotient zweier Größen \emph{derselben} Dimension; keine Einheitenwahl kürzt
    ihn weg. Solche Verhältnisse sind, anders als $\alpha$, nicht durch eine
    Normierung erreichbar -- $\alpha$ verbirgt eine Ladungseinheit, ein
    Massenverhältnis verbirgt nichts.
\end{itemize}

\textbf{[SETZUNG]} Damit steht die Trennlinie: setzbar auf Eins $\Rightarrow$
Konvention; nicht setzbar $\Rightarrow$ Inhalt. $\alpha$ steht trotz seiner
Dimensionslosigkeit auf der Konventionsseite, weil es eine Ladungseinheit trägt;
$m_\mu/m_e$ steht auf der Inhaltsseite. Der einzige verbleibende Inhalt der
Theorie ist $\xi$ und eine einzige absolute Skala.

\section{Zwei Seiten einer Medaille -- und trotzdem nur eine grundlegend}

\textbf{[K]} Hier droht ein Widerspruch, den man ausräumen muss. A130 nennt
$\alpha$ und $\xi$ \enquote{zwei Seiten einer Medaille}, einen einzigen
Freiheitsgrad, verbunden durch
\[
  \alpha_\text{SI} = \xi\cdot E_0^2 .
\]
Dieses Dokument sagt zugleich, $\alpha$ sei nicht grundlegend, $\xi$ dagegen
sehr wohl. Wie passt beides zusammen?

\textbf{[K]} Es passt, sobald man \enquote{ein Freiheitsgrad} richtig liest: die
Aussage betrifft den \emph{Informationsgehalt}, nicht den \emph{Rang}. $\alpha_\text{SI}$
und $\xi$ tragen \emph{dieselbe} geometrische Information -- wer $\xi$ kennt und
die Skala $E_0$, kennt $\alpha_\text{SI}$, und umgekehrt. In diesem Sinn ist es
eine Größe, nicht zwei. Das ist die Medaille.

\textbf{[B]} Die beiden Seiten sind aber nicht gleichrangig. $\xi$ ist diese
Information \emph{blank}: ein reines Verhältnis, in jedem Maßsystem $4/30000$.
$\alpha_\text{SI}$ ist \emph{dieselbe} Information \emph{eingekleidet}:
$\alpha_\text{SI} = \xi\cdot E_0^2$, wobei der Faktor $E_0^2$ reine Skala ist --
die Einheiten-Kleidung. Der berühmte Wert $1/137$ ist also nicht $\xi$ neben
$\alpha$, sondern $\xi$ \emph{im Kostüm} von $E_0^2$.

\textbf{[K]} Die Diagnose dieses Dokuments bestätigt genau das. Zieht man die
Einheiten ab -- setzt man $\alpha = 1$ in natürlichen Einheiten --, trägt
$\alpha$ nichts mehr: die Information ist nicht verschwunden, sie steckt
vollständig in $\xi$ und der einen Skala. $\xi$ übersteht jeden Einheitenwechsel;
$\alpha = 1$ ist danach leer. Die Medaille bleibt also erhalten, aber sie
\emph{ist} $\xi$: $\alpha_\text{SI}$ ist, wie diese Münze im SI-Licht aussieht,
nicht eine zweite Münze.

\textbf{[SETZUNG]} Kurz: \enquote{zwei Seiten} heißt gleicher Gehalt, nicht
gleicher Rang. Ein Freiheitsgrad -- und seine grundlegende Gestalt ist $\xi$, die
eingekleidete $\alpha_\text{SI}$. Der scheinbare Widerspruch entsteht nur, wenn
man \enquote{äquivalent} mit \enquote{gleich fundamental} verwechselt.

\section{Wozu dann überhaupt die SI-Rechnung?}

\textbf{[K]} Wenn in natürlichen Einheiten alles außer $\xi$ und einer Skala auf
Eins fällt, drängt sich die Frage auf, wozu man die aufwendigeren
SI-Rechnungen, die Brückenkonstanten und die mehrfachen Herleitungswege
überhaupt braucht. Die Antwort trennt drei Dinge, die leicht verwechselt werden.

\textbf{[K]} \emph{Erstens: Übersetzung.} Gemessen wird im Labor, und Labore
rechnen in SI. Die SI-Kleidung ist die \emph{Übersetzung} des geometrischen
Kerns in ablesbare Zahlen -- unverzichtbar, aber nicht selbst Physik. Man
braucht sie, weil das Messgerät in Metern und Sekunden spricht, nicht weil der
Inhalt sie verlangt. Ein Ergebnis in natürlichen Einheiten ist vollständig; die
SI-Zahl ist seine Aussprache in einer bestimmten Sprache.

\textbf{[B]} \emph{Zweitens: Überbestimmung als Test.} Die mehrfachen
Rechenwege fügen dem Inhalt nichts hinzu -- sie \emph{prüfen} ihn. Wenn zwei
Wege zur selben Energieskala führen und ihre Differenz genau der abgeleitete
Fraktalfaktor ist (A130), oder wenn drei Herleitungen der Leptonmassen
denselben Wert treffen (A110), dann ist das kein zusätzlicher Freiheitsgrad,
sondern eine Kreuzkontrolle. Ein einzelner Weg könnte ein verkleideter Fit
sein; mehrere unabhängige Wege, die zusammenfallen, können es nicht sein. Die
Redundanz ist die Härtung, nicht der Inhalt.

\textbf{[K]} \emph{Drittens: Anschlussfähigkeit.} Verschiedene Wege binden den
einen $\xi$-Inhalt an verschiedene bekannte Physik -- Yukawa-Kopplung,
Koide-Relation, RG-Lauf. Das macht denselben Kern aus mehreren vertrauten
Richtungen prüfbar und verhindert, dass er nur in seiner eigenen Sprache
stimmig aussieht. Der Inhalt bleibt einer; die Zugänge sind viele, damit er von
außen kontrollierbar ist.

\textbf{[SETZUNG]} Kurz: die SI-Rechnung ist Übersetzung, die mehrfachen Wege
sind Kontrolle. Beide sind nötig, um den geometrischen Kern mit der Messung zu
verbinden und ihn zu härten -- keiner von beiden fügt dem Kern einen
Freiheitsgrad hinzu. Der Kern bleibt $\xi$ und eine Skala.

\section{Prüfskript}

\begin{itemize}
  \item Die Diagnose maschinell: für $c, \hbar, k_B, \alpha$ wird gezeigt, dass
    eine Einheiten-/Normierungswahl sie auf Eins bringt (also Konvention); für
    $\xi$ und $m_\mu/m_e$, dass keine Einheitenwahl das leistet (also Inhalt);
    und die Kreuzkontrolle, dass zwei Wege denselben Wert treffen, ohne einen
    Freiheitsgrad hinzuzufügen:
    \texttt{python/A\_Serie\_Skripte/a135\_eins\_diagnose.py}
\end{itemize}

\section{Was offen bleibt}

\textbf{Erstens bleibt die eine Skala eine Eingabe.} Der Satz reduziert alles auf
$\xi$ und eine absolute Skala; welchen Wert diese Skala hat, ist selbst nicht
hergeleitet (A130), sondern der eine verbleibende Kalibrierungspunkt. Die Skala
ist \emph{nicht} auf Eins setzbar in dem Sinn, dass ihr Wert Inhalt trägt --
aber ihr Zahlenwert ist die eine Eingabe, die das System braucht.

\textbf{Zweitens ist die Diagnose kein Beweis der Vollständigkeit.} Dass alles
\emph{Bisherige} entweder auf Eins fällt oder $\xi$ ist, zeigt die
Sparsamkeit der Theorie; es beweist nicht, dass keine weitere unabhängige Größe
je nötig wird. Der Satz ordnet das Vorhandene, er schließt Künftiges nicht aus.

\section{Ersetzt}

Dieses Dokument nimmt auf: Dok.~043 und Dok.~044 (Auflösung der Konstanten),
Dok.~101 (Zirkularität der Konstanten, ein Freiheitsgrad). Es zieht die
Diagnose zusammen, die in A085 (Dualität) und A130 ($\alpha$) je für einen Fall
geführt wird, und beantwortet die Frage nach dem Zweck der SI-Rechnung.
Eingearbeitet: P40.

\section{Verwendung von KI-Werkzeugen}

Dieses Dokument ist unter Verwendung KI-gestützter Werkzeuge entstanden. Die
Fragestellungen, die Setzungen und alle inhaltlichen Entscheidungen stammen vom
Autor; die Werkzeuge formulieren aus, rechnen nach und erstellen Prüfskripte.
Die Verantwortung für Inhalt und Fehler liegt vollständig beim Autor. Der
ausführliche Wortlaut steht in A010.
